% ============================================================
% Question Bank: ch03-07 Converting Rational Numbers to Decimals
% Topic Title: Converting Rational Numbers to Decimals
% CCSS: 7.NS.A.2d
% Questions: 27 (15 MC + 6 SA + 6 Graphical)
% ============================================================

% --- Multiple Choice Questions (15) ---

\begin{practiceQuestion}{3-7-q01}{mc}
\begin{questionText}
What is $\dfrac{3}{8}$ written as a decimal?
\end{questionText}
\choiceA{$0.35$}
\choiceB{$0.375$}
\choiceC{$0.38$}
\choiceD{$0.3$}
\correctAnswer{B}
\explanation{Divide $3 \div 8 = 0.375$. This is a terminating decimal because $8 = 2^3$ has only factors of $2$.}
\end{practiceQuestion}

\begin{practiceQuestion}{3-7-q02}{mc}
\begin{questionText}
What is $\dfrac{5}{9}$ written as a decimal?
\end{questionText}
\choiceA{$0.5$}
\choiceB{$0.55$}
\choiceC{$0.\overline{5}$}
\choiceD{$0.59$}
\correctAnswer{C}
\explanation{$5 \div 9 = 0.555\ldots = 0.\overline{5}$. The digit $5$ repeats because $9$ has a factor of $3$ (not $2$ or $5$).}
\end{practiceQuestion}

\begin{practiceQuestion}{3-7-q03}{mc}
\begin{questionText}
Which fraction converts to a terminating decimal?
\end{questionText}
\choiceA{$\dfrac{1}{3}$}
\choiceB{$\dfrac{7}{12}$}
\choiceC{$\dfrac{3}{16}$}
\choiceD{$\dfrac{5}{6}$}
\correctAnswer{C}
\explanation{A terminating decimal has a denominator whose only prime factors are $2$ and/or $5$. $16 = 2^4$ qualifies. The others have factor $3$.}
\end{practiceQuestion}

\begin{practiceQuestion}{3-7-q04}{mc}
\begin{questionText}
What is $-\dfrac{7}{11}$ written as a decimal?
\end{questionText}
\choiceA{$-0.\overline{63}$}
\choiceB{$-0.63$}
\choiceC{$-0.636$}
\choiceD{$-0.6\overline{3}$}
\correctAnswer{A}
\explanation{$7 \div 11 = 0.636363\ldots = 0.\overline{63}$. The pair $63$ repeats. The negative sign carries over: $-0.\overline{63}$.}
\end{practiceQuestion}

\begin{practiceQuestion}{3-7-q05}{mc}
\begin{questionText}
Which decimal is equivalent to $\dfrac{2}{5}$?
\end{questionText}
\choiceA{$0.25$}
\choiceB{$0.4$}
\choiceC{$0.2$}
\choiceD{$0.45$}
\correctAnswer{B}
\explanation{$2 \div 5 = 0.4$. Or multiply: $\dfrac{2}{5} \times \dfrac{2}{2} = \dfrac{4}{10} = 0.4$.}
\end{practiceQuestion}

\begin{practiceQuestion}{3-7-q06}{mc}
\begin{questionText}
What is $0.\overline{3}$ written as a fraction?
\end{questionText}
\choiceA{$\dfrac{3}{10}$}
\choiceB{$\dfrac{3}{9}$}
\choiceC{$\dfrac{1}{3}$}
\choiceD{Both B and C}
\correctAnswer{D}
\explanation{$0.\overline{3} = 0.333\ldots = \dfrac{3}{9} = \dfrac{1}{3}$. Both B and C represent the same value.}
\end{practiceQuestion}

\begin{practiceQuestion}{3-7-q07}{mc}
\begin{questionText}
What is $\dfrac{11}{4}$ written as a decimal?
\end{questionText}
\choiceA{$2.5$}
\choiceB{$2.75$}
\choiceC{$2.25$}
\choiceD{$2.4$}
\correctAnswer{B}
\explanation{$11 \div 4 = 2.75$. This terminates because $4 = 2^2$.}
\end{practiceQuestion}

\begin{practiceQuestion}{3-7-q08}{mc}
\begin{questionText}
Which decimal represents $\dfrac{1}{6}$?
\end{questionText}
\choiceA{$0.1\overline{6}$}
\choiceB{$0.\overline{16}$}
\choiceC{$0.16$}
\choiceD{$0.166$}
\correctAnswer{A}
\explanation{$1 \div 6 = 0.1666\ldots = 0.1\overline{6}$. Only the $6$ repeats; the $1$ appears once.}
\end{practiceQuestion}

\begin{practiceQuestion}{3-7-q09}{mc}
\begin{questionText}
Which statement about $\dfrac{7}{20}$ is true?
\end{questionText}
\choiceA{It is a repeating decimal}
\choiceB{It equals $0.35$}
\choiceC{It equals $0.37$}
\choiceD{It is irrational}
\correctAnswer{B}
\explanation{$20 = 2^2 \times 5$, so it terminates. $7 \div 20 = 0.35$.}
\end{practiceQuestion}

\begin{practiceQuestion}{3-7-q10}{mc}
\begin{questionText}
What type of decimal does $\dfrac{4}{7}$ produce?
\end{questionText}
\choiceA{Terminating}
\choiceB{Repeating}
\choiceC{Neither — it is irrational}
\choiceD{Both terminating and repeating}
\correctAnswer{B}
\explanation{$7$ has a prime factor other than $2$ or $5$, so $\dfrac{4}{7}$ is a repeating decimal: $0.\overline{571428}$.}
\end{practiceQuestion}

\begin{practiceQuestion}{3-7-q11}{mc}
\begin{questionText}
What is $-\dfrac{3}{20}$ as a decimal?
\end{questionText}
\choiceA{$-0.15$}
\choiceB{$-0.35$}
\choiceC{$-0.2$}
\choiceD{$-0.03$}
\correctAnswer{A}
\explanation{$3 \div 20 = 0.15$. The negative sign gives $-0.15$.}
\end{practiceQuestion}

\begin{practiceQuestion}{3-7-q12}{mc}
\begin{questionText}
Which fraction is equivalent to $0.625$?
\end{questionText}
\choiceA{$\dfrac{5}{8}$}
\choiceB{$\dfrac{6}{8}$}
\choiceC{$\dfrac{5}{6}$}
\choiceD{$\dfrac{3}{5}$}
\correctAnswer{A}
\explanation{$0.625 = \dfrac{625}{1000} = \dfrac{5}{8}$. Or check: $5 \div 8 = 0.625$.}
\end{practiceQuestion}

\begin{practiceQuestion}{3-7-q13}{mc}
\begin{questionText}
What bar notation represents $0.272727\ldots$?
\end{questionText}
\choiceA{$0.\overline{2}$}
\choiceB{$0.2\overline{7}$}
\choiceC{$0.\overline{27}$}
\choiceD{$0.\overline{272}$}
\correctAnswer{C}
\explanation{The digits $27$ repeat as a pair: $0.272727\ldots = 0.\overline{27}$.}
\end{practiceQuestion}

\begin{practiceQuestion}{3-7-q14}{mc}
\begin{questionText}
Between which two integers does $-\dfrac{13}{4}$ lie?
\end{questionText}
\choiceA{$-3$ and $-2$}
\choiceB{$-4$ and $-3$}
\choiceC{$-5$ and $-4$}
\choiceD{$3$ and $4$}
\correctAnswer{B}
\explanation{$-\dfrac{13}{4} = -3.25$, which lies between $-4$ and $-3$.}
\end{practiceQuestion}

\begin{practiceQuestion}{3-7-q15}{mc}
\begin{questionText}
Which list is in order from least to greatest?
\end{questionText}
\choiceA{$\dfrac{1}{3},\ 0.3,\ \dfrac{3}{10}$}
\choiceB{$\dfrac{3}{10},\ 0.3,\ \dfrac{1}{3}$}
\choiceC{$0.3,\ \dfrac{3}{10},\ \dfrac{1}{3}$}
\choiceD{$\dfrac{3}{10},\ \dfrac{1}{3},\ 0.3$}
\correctAnswer{C}
\explanation{$0.3 = 0.3$, $\dfrac{3}{10} = 0.3$, $\dfrac{1}{3} = 0.\overline{3} \approx 0.333$. Since $0.3 = \dfrac{3}{10} < \dfrac{1}{3}$, C is correct (both $0.3$ and $\dfrac{3}{10}$ tie, then $\dfrac{1}{3}$ is greatest).}
\end{practiceQuestion}

% --- Short Answer Questions (6) ---

\begin{practiceQuestion}{3-7-q16}{sa}
\begin{questionText}
Convert $\dfrac{5}{12}$ to a decimal. Use bar notation for repeating digits.
\end{questionText}
\correctAnswer{$0.41\overline{6}$}
\explanation{$5 \div 12 = 0.41666\ldots = 0.41\overline{6}$. The $6$ repeats.}
\end{practiceQuestion}

\begin{practiceQuestion}{3-7-q17}{sa}
\begin{questionText}
Convert $\dfrac{7}{25}$ to a decimal.
\end{questionText}
\correctAnswer{$0.28$}
\explanation{$7 \div 25 = 0.28$. Since $25 = 5^2$, this is a terminating decimal.}
\end{practiceQuestion}

\begin{practiceQuestion}{3-7-q18}{sa}
\begin{questionText}
Is $\dfrac{9}{15}$ a terminating or repeating decimal? Explain.
\end{questionText}
\correctAnswer{Repeating: $0.\overline{6}$}
\explanation{Simplify: $\dfrac{9}{15} = \dfrac{3}{5} = 0.6$. Wait — $\dfrac{3}{5}$ terminates because $5$ is a factor. So $\dfrac{9}{15} = 0.6$ terminates.}
\end{practiceQuestion}

\begin{practiceQuestion}{3-7-q19}{sa}
\begin{questionText}
What fraction in simplest form equals $0.\overline{45}$?
\end{questionText}
\correctAnswer{$\dfrac{5}{11}$}
\explanation{Let $x = 0.\overline{45}$. Then $100x = 45.\overline{45}$. So $99x = 45$, giving $x = \dfrac{45}{99} = \dfrac{5}{11}$.}
\end{practiceQuestion}

\begin{practiceQuestion}{3-7-q20}{sa}
\begin{questionText}
Convert $-3\dfrac{5}{8}$ to a decimal.
\end{questionText}
\correctAnswer{$-3.625$}
\explanation{$\dfrac{5}{8} = 0.625$. So $-3\dfrac{5}{8} = -3.625$.}
\end{practiceQuestion}

\begin{practiceQuestion}{3-7-q21}{sa}
\begin{questionText}
Place $\dfrac{7}{9}$ on the number line between which two tenths?
\end{questionText}
\correctAnswer{$0.7$ and $0.8$}
\explanation{$\dfrac{7}{9} = 0.\overline{7} \approx 0.778$. This value lies between $0.7$ and $0.8$.}
\end{practiceQuestion}

% --- Graphical Questions (3 GMC + 3 GSA) ---

\begin{practiceQuestion}{3-7-q22}{gmc}
\begin{questionText}
The number line shows point $P$. Which fraction is closest to $P$?

\begin{center}
\begin{tikzpicture}
  \draw[thick,->] (-0.5,0) -- (2.5,0);
  \draw (0,0.1) -- (0,-0.1) node[below,font=\small] {$0$};
  \draw (2,0.1) -- (2,-0.1) node[below,font=\small] {$1$};
  \foreach \x in {0.2,0.4,...,1.8} \draw (\x,0.05) -- (\x,-0.05);
  \fill[funBlue] (0.7,0) circle (0.1) node[above=4pt,font=\small] {$P$};
\end{tikzpicture}
\end{center}
\end{questionText}
\choiceA{$\dfrac{1}{4}$}
\choiceB{$\dfrac{1}{3}$}
\choiceC{$\dfrac{3}{8}$}
\choiceD{$\dfrac{7}{20}$}
\correctAnswer{D}
\explanation{$P$ is at approximately $0.35$. $\dfrac{1}{4} = 0.25$, $\dfrac{1}{3} \approx 0.33$, $\dfrac{3}{8} = 0.375$, $\dfrac{7}{20} = 0.35$. The closest is $\dfrac{7}{20} = 0.35$.}
\end{practiceQuestion}

\begin{practiceQuestion}{3-7-q23}{gmc}
\begin{questionText}
The table shows fractions and their decimal forms. Which row contains an error?

\begin{center}
\begin{tabular}{|c|c|c|}
\hline
\textbf{Row} & \textbf{Fraction} & \textbf{Decimal} \\
\hline
$1$ & $\dfrac{3}{4}$ & $0.75$ \\ \hline
$2$ & $\dfrac{2}{3}$ & $0.6\overline{6}$ \\ \hline
$3$ & $\dfrac{5}{16}$ & $0.3125$ \\ \hline
$4$ & $\dfrac{4}{9}$ & $0.4$ \\ \hline
\end{tabular}
\end{center}
\end{questionText}
\choiceA{Row 1}
\choiceB{Row 2}
\choiceC{Row 3}
\choiceD{Row 4}
\correctAnswer{D}
\explanation{$\dfrac{4}{9} = 0.\overline{4}$, not $0.4$. The bar over the $4$ is missing; it is a repeating decimal.}
\end{practiceQuestion}

\begin{practiceQuestion}{3-7-q24}{gmc}
\begin{questionText}
The long division below shows $7 \div 11$. What digit fills the blank to continue the pattern?

\begin{center}
\begin{tabular}{l}
$0.6363\ldots$ \\
\hline
$11 \overline{)\ 7.0000}$ \\
$\phantom{1}66$ \\
\hline
$\phantom{11}40$ \\
$\phantom{11}33$ \\
\hline
$\phantom{111}70$ \\
$\phantom{111}66$ \\
\hline
$\phantom{1111}\boxed{?}0$ \\
\end{tabular}
\end{center}
\end{questionText}
\choiceA{$3$}
\choiceB{$4$}
\choiceC{$6$}
\choiceD{$7$}
\correctAnswer{B}
\explanation{$70 - 66 = 4$, so the remainder is $4$ and the next step uses $40$. The repeating cycle continues: $\ldots63\ldots$}
\end{practiceQuestion}

\begin{practiceQuestion}{3-7-q25}{gsa}
\begin{questionText}
Convert $\dfrac{5}{6}$ to a decimal using long division. Write your answer using bar notation.

\begin{center}
\begin{tabular}{l}
$0.8333\ldots$ \\
\hline
$6 \overline{)\ 5.0000}$ \\
$\phantom{1}48$ \\
\hline
$\phantom{11}20$ \\
$\phantom{11}18$ \\
\hline
$\phantom{111}20$ \\
$\phantom{111}18$ \\
\hline
$\phantom{1111}20$ \\
\end{tabular}
\end{center}
\end{questionText}
\correctAnswer{$0.8\overline{3}$}
\explanation{$5 \div 6 = 0.8333\ldots = 0.8\overline{3}$. The $3$ repeats indefinitely after the initial $8$.}
\end{practiceQuestion}

\begin{practiceQuestion}{3-7-q26}{gsa}
\begin{questionText}
The number line shows points $A$ and $B$. Convert each to a fraction in simplest form and find $A + B$.

\begin{center}
\begin{tikzpicture}
  \draw[thick,->] (-0.5,0) -- (3.5,0);
  \draw (0,0.1) -- (0,-0.1) node[below,font=\small] {$0$};
  \draw (1,0.1) -- (1,-0.1) node[below,font=\small] {$0.5$};
  \draw (2,0.1) -- (2,-0.1) node[below,font=\small] {$1.0$};
  \draw (3,0.1) -- (3,-0.1) node[below,font=\small] {$1.5$};
  \fill[funBlue] (0.5,0) circle (0.1) node[above=4pt,font=\small] {$A$};
  \fill[funRed] (2.5,0) circle (0.1) node[above=4pt,font=\small] {$B$};
\end{tikzpicture}
\end{center}
\end{questionText}
\correctAnswer{$\dfrac{7}{4}$ or $1.75$ or $1\dfrac{3}{4}$}
\explanation{$A = 0.25 = \dfrac{1}{4}$, $B = 1.25 = \dfrac{5}{4}$. Sum: $\dfrac{1}{4} + \dfrac{5}{4} = \dfrac{6}{4} = \dfrac{3}{2} = 1.5$. Actually $A = 0.25$ and $B = 1.25$, so sum $= 1.5$.}
\end{practiceQuestion}

\begin{practiceQuestion}{3-7-q27}{gsa}
\begin{questionText}
Complete the conversion table. Fill in the missing decimal value.

\begin{center}
\begin{tabular}{|c|c|}
\hline
\textbf{Fraction} & \textbf{Decimal} \\
\hline
$\dfrac{1}{8}$ & $0.125$ \\ \hline
$\dfrac{3}{8}$ & $0.375$ \\ \hline
$\dfrac{5}{8}$ & ? \\ \hline
$\dfrac{7}{8}$ & $0.875$ \\ \hline
\end{tabular}
\end{center}
\end{questionText}
\correctAnswer{$0.625$}
\explanation{$5 \div 8 = 0.625$. The pattern: each value increases by $0.250$ ($\dfrac{2}{8}$).}
\end{practiceQuestion}
